% Options for packages loaded elsewhere
\PassOptionsToPackage{unicode}{hyperref}
\PassOptionsToPackage{hyphens}{url}
\PassOptionsToPackage{dvipsnames,svgnames,x11names}{xcolor}
\documentclass[
  spanish,
  10pt,
  a4paper,
]{article}
\usepackage{xcolor}
\usepackage[top=2.2cm,bottom=2.2cm,left=2.1cm,right=2.1cm]{geometry}
\usepackage{amsmath,amssymb}
\setcounter{secnumdepth}{5}
\usepackage{iftex}
\ifPDFTeX
  \usepackage[T1]{fontenc}
  \usepackage[utf8]{inputenc}
  \usepackage{textcomp} % provide euro and other symbols
\else % if luatex or xetex
  \usepackage{unicode-math} % this also loads fontspec
  \defaultfontfeatures{Scale=MatchLowercase}
  \defaultfontfeatures[\rmfamily]{Ligatures=TeX,Scale=1}
\fi
\usepackage{lmodern}
\ifPDFTeX\else
  % xetex/luatex font selection
\fi
% Use upquote if available, for straight quotes in verbatim environments
\IfFileExists{upquote.sty}{\usepackage{upquote}}{}
\IfFileExists{microtype.sty}{% use microtype if available
  \usepackage[]{microtype}
  \UseMicrotypeSet[protrusion]{basicmath} % disable protrusion for tt fonts
}{}
\makeatletter
\@ifundefined{KOMAClassName}{% if non-KOMA class
  \IfFileExists{parskip.sty}{%
    \usepackage{parskip}
  }{% else
    \setlength{\parindent}{0pt}
    \setlength{\parskip}{6pt plus 2pt minus 1pt}}
}{% if KOMA class
  \KOMAoptions{parskip=half}}
\makeatother
\ifLuaTeX
\usepackage[bidi=basic,shorthands=off]{babel}
\else
\usepackage[bidi=default,shorthands=off]{babel}
\fi
\ifLuaTeX
  \usepackage{selnolig} % disable illegal ligatures
\fi
\setlength{\emergencystretch}{3em} % prevent overfull lines
\providecommand{\tightlist}{%
  \setlength{\itemsep}{0pt}\setlength{\parskip}{0pt}}
% Shared visual layer for the generated Kaleido FlowTwin reports.
\usepackage{helvet}
\renewcommand{\familydefault}{\sfdefault}
\usepackage{microtype}
\usepackage{xcolor}
\usepackage{titlesec}
\usepackage{fancyhdr}
\usepackage{enumitem}
\usepackage{booktabs}
\usepackage{xurl}
\usepackage{lastpage}
\usepackage{etoolbox}

\definecolor{KaleidoNavy}{HTML}{0A2342}
\definecolor{KaleidoBlue}{HTML}{176B87}
\definecolor{KaleidoCyan}{HTML}{3BB7C4}
\definecolor{KaleidoMint}{HTML}{DDF5F1}
\definecolor{KaleidoFog}{HTML}{F2F5F7}
\definecolor{KaleidoSlate}{HTML}{5D6A7A}
\definecolor{KaleidoCoral}{HTML}{F26B5B}

\titleformat{\section}
  {\Large\bfseries\color{KaleidoNavy}}
  {\thesection}{0.75em}{}[\vspace{0.2em}\color{KaleidoCyan}\titlerule]
\titleformat{\subsection}
  {\large\bfseries\color{KaleidoBlue}}
  {\thesubsection}{0.65em}{}
\titleformat{\subsubsection}
  {\normalsize\bfseries\color{KaleidoNavy}}
  {\thesubsubsection}{0.55em}{}
\titlespacing*{\section}{0pt}{2.2em}{1em}
\titlespacing*{\subsection}{0pt}{1.6em}{0.65em}

\pagestyle{fancy}
\fancyhf{}
\fancyhead[L]{\small\bfseries\color{KaleidoNavy}KALEIDO FLOWTWIN}
\fancyhead[R]{\small\color{KaleidoSlate}DOSSIER TECNICO}
\fancyfoot[L]{\scriptsize\color{KaleidoSlate}EVOCON Solutions GmbH}
\fancyfoot[R]{\scriptsize\color{KaleidoSlate}\thepage/\pageref*{LastPage}}
\renewcommand{\headrulewidth}{0.6pt}
\renewcommand{\headrule}{\hbox to\headwidth{\color{KaleidoCyan}\leaders\hrule height \headrulewidth\hfill}}
\renewcommand{\footrulewidth}{0pt}
\setlength{\headheight}{20pt}

\setlist[itemize]{leftmargin=1.4em,itemsep=0.25em,topsep=0.35em}
\setlist[enumerate]{leftmargin=1.7em,itemsep=0.25em,topsep=0.35em}
\setlength{\parindent}{0pt}
\setlength{\parskip}{0.55em}
\setlength{\emergencystretch}{3em}
\renewcommand{\arraystretch}{1.18}

\AtBeginEnvironment{quote}{%
  \color{KaleidoNavy}\itshape
  \setlength{\leftskip}{1em}
  \setlength{\rightskip}{1em}
}

\makeatletter
\renewcommand{\maketitle}{%
  \hypersetup{pageanchor=false}
  \begin{titlepage}
    \pagecolor{KaleidoNavy}
    \color{white}
    \vspace*{1.4cm}
    {\color{KaleidoCyan}\rule{2.4cm}{2.5pt}\par}
    \vspace{1.1cm}
    {\Huge\bfseries\@title\par}
    \vspace{0.65cm}
    {\Large\color{KaleidoMint}Kaleido FlowTwin\par}
    \vfill
    {\large Álvaro Schwiedop Souto\par}
    \vspace{0.15cm}
    {\normalsize\color{KaleidoMint}Industrial AI \& Predictive Quality Lead\par}
    \vspace{0.15cm}
    {\normalsize\color{KaleidoMint}EVOCON Solutions GmbH\par}
    \vspace{0.25cm}
    {\normalsize\color{KaleidoMint}\@date\par}
  \end{titlepage}
  \hypersetup{pageanchor=true}
  \nopagecolor
  \color{black}
}
\makeatother
\usepackage{bookmark}
\IfFileExists{xurl.sty}{\usepackage{xurl}}{} % add URL line breaks if available
\urlstyle{same}
% fallback for those not using the hyperref driver hyperxmp:
\makeatletter
\@ifundefined{xmpquote}{\newcommand{\xmpquote}[1]{#1}}{}
\makeatother
\hypersetup{
  pdftitle={Kaleido: investigacion de encaje},
  pdflang={es},
  colorlinks=true,
  linkcolor={KaleidoBlue},
  filecolor={Maroon},
  citecolor={Blue},
  urlcolor={KaleidoBlue},
  pdfcreator={LaTeX via pandoc}}

\title{Kaleido: investigacion de encaje}
\author{}
\date{15 julio 2026}

\begin{document}
\maketitle

{
\setcounter{tocdepth}{3}
\tableofcontents
}
\section{Kaleido: investigacion de
encaje}\label{kaleido-investigacion-de-encaje}

Fecha de corte: 15-07-2026. Se han priorizado paginas oficiales de
Kaleido.

\subsection{Lo que hace Kaleido}\label{lo-que-hace-kaleido}

Kaleido Ideas \& Logistics opera soluciones integradas por mar, tierra y
aire y combina transporte, aduanas, almacenaje, trincaje, operaciones
portuarias, ingenieria e innovacion. Sus proyectos muestran carga de
proyecto y breakbulk de alta complejidad: granito, estructuras eolicas,
material ferroviario, equipos industriales, buques y cargas
sobredimensionadas.

La terminal de Vigo declarada por Kaleido tiene 50.000 m2, almacen
interior y exterior, zonas refrigeradas y servicios de carga, descarga,
manipulacion, agencia, trincaje, soldadura, aduanas, ingenieria,
transporte y chartering.

Fuentes:

\begin{itemize}
\tightlist
\item
  \url{https://www.kaleidologistics.com/en/}
\item
  \url{https://www.kaleidologistics.com/en/port-operations/}
\item
  \url{https://www.kaleidologistics.com/en/engineering-solutions/}
\end{itemize}

\subsection{Lo que ya ha construido Kaleido
Tech}\label{lo-que-ya-ha-construido-kaleido-tech}

\subsubsection{Trace Port}\label{trace-port}

Plataforma web y movil para registrar, analizar y coordinar operaciones
portuarias. Incluye proyectos, packing lists, turnos, equipos, eventos,
fotografias, danos/incidencias, informes, timestamps, historial y API.

Consecuencia para FlowTwin: es el mejor punto de entrada. Ya instrumenta
la traza operacional que necesita un predictor y el valor se puede
devolver dentro de la misma experiencia de usuario.

Fuente:
\url{https://www.kaleidologistics.com/en/productos-ktech/trace-port/}

\subsubsection{Freight Intelligence}\label{freight-intelligence}

Centraliza seguimiento de contenedores, integra mas de 180 navieras
segun la pagina del producto, aporta posicion y estado 24/7, alertas,
informes e integracion con sistemas de gestion.

Consecuencia: es una segunda vertical para excepciones/ETA, pero compite
en un mercado mas poblado y gran parte de la senal depende de terceros.
Conviene abordarla despues de demostrar el patron en Trace Port.

Fuente:
\url{https://www.kaleidologistics.com/en/productos-ktech/freight-intellligence/}

\subsubsection{Shipping Board}\label{shipping-board}

Coordina expediciones y recepciones entre almacen, transporte y
logistica, digitaliza documentos, automatiza avisos y exporta
estadisticas.

Consecuencia: ofrece secuencias de eventos utiles para predecir tiempo
restante, atascos y retrasos de recepcion/expedicion.

Fuente:
\url{https://www.kaleidologistics.com/en/productos-ktech/shipping-board/}

\subsubsection{Karbon Track}\label{karbon-track}

Calcula y reporta CO2e por modo y etapa, con importacion Excel/CSV/API.
Es un producto de reporting y cumplimiento. FlowTwin podria usar sus
costes de emision como objetivo secundario de escenarios, no como primer
caso.

Fuente:
\url{https://www.kaleidologistics.com/en/productos-ktech/karbon-track/}

\subsection{Lectura del correo
recibido}\label{lectura-del-correo-recibido}

Manuel Alonso Montenegro indica que Kaleido no cuenta con un historico
de operativas muy rico. Esto no es un rechazo; es el principal requisito
de diseno.

Interpretacion tecnica:

\begin{itemize}
\tightlist
\item
  no hay base para prometer un gran modelo supervisado;
\item
  el primer entregable debe medir riqueza, cobertura y trazabilidad;
\item
  una unidad de negocio con herramientas propias puede empezar a generar
  mejor dato de forma prospectiva;
\item
  los eventos sin etiqueta siguen siendo utiles para process mining,
  estimacion descriptiva y pretraining autosupervisado;
\item
  un predictor solo se promueve si mejora baselines simples bajo
  holdout.
\end{itemize}

Interpretacion comercial:

\begin{itemize}
\tightlist
\item
  Kaleido conoce el problema y valora el encaje, pero no comprara
  investigacion abstracta;
\item
  necesita un piloto que pueda terminar honestamente en tres resultados:
  \texttt{go}, \texttt{instrument\ first} o \texttt{no\ signal};
\item
  la propuesta debe mejorar su catalogo y generar una nueva linea
  vendible, no duplicar una funcionalidad que ya posee.
\end{itemize}

\subsection{Caso inicial recomendado}\label{caso-inicial-recomendado}

\subsubsection{Riesgo de desviacion de operacion/turno en Trace
Port}\label{riesgo-de-desviacion-de-operacionturno-en-trace-port}

Unidad: una operacion o turno de carga, descarga o manipulacion.

Pregunta:

\begin{quote}
Con el prefijo de eventos observado hasta ahora, el plan y el contexto,
¿se terminara dentro del plan y que probabilidad existe de una
desviacion material en las proximas 2/4/8 horas?
\end{quote}

Subsalidas:

\begin{itemize}
\tightlist
\item
  tiempo restante P50/P90;
\item
  milestone siguiente y hora probable;
\item
  probabilidad de exceder plan;
\item
  cuello de botella actual;
\item
  factor de riesgo: recurso, secuencia, espera, incidencia, clima o
  documentacion;
\item
  recomendacion dentro de una lista validada por Kaleido.
\end{itemize}

Por que este caso:

\begin{itemize}
\tightlist
\item
  utiliza el dato que Trace Port ya declara capturar;
\item
  el resultado es comprensible por operaciones;
\item
  se puede evaluar con duracion, lead time y falsas alarmas;
\item
  puede reducir sobrecoste, horas improductivas y coordinacion manual;
\item
  se empaqueta como modulo premium del producto existente;
\item
  el mismo esquema se transfiere luego a Shipping Board.
\end{itemize}

\subsection{Opciones descartadas como primer
piloto}\label{opciones-descartadas-como-primer-piloto}

\subsubsection{Un world model generativo de
video}\label{un-world-model-generativo-de-video}

Demasiado dato y computo, sin relacion directa con el cuello de botella
manifestado. Generar imagenes plausibles no demuestra fidelidad
operacional.

\subsubsection{Gemelo 3D LiDAR del
puerto}\label{gemelo-3d-lidar-del-puerto}

Interesante para inventario, seguridad espacial y simulacion de
maniobras, pero requiere sensores, mapas y un problema 3D concreto. No
resuelve el cold start de eventos y debe ser una fase futura.

\subsubsection{Robotica autonoma}\label{robotica-autonoma}

SoftNav-JEPA aporta ideas, pero el stack actual es un proxy cinematico y
sus resultados pre-fix estan invalidados. No es un producto listo para
puerto.

\subsubsection{ETA de contenedores como primer
caso}\label{eta-de-contenedores-como-primer-caso}

Factible gracias a Freight Intelligence, pero menos diferencial que
explotar el conocimiento de operacion portuaria y de carga de proyecto
de Kaleido.

\subsection{Como gana dinero Kaleido}\label{como-gana-dinero-kaleido}

Hipotesis comerciales a validar con Manuel:

\begin{enumerate}
\def\labelenumi{\arabic{enumi}.}
\tightlist
\item
  modulo SaaS premium de \texttt{predictive\ operations} por
  terminal/proyecto;
\item
  add-on white-label para clientes de Trace Port;
\item
  servicio de implantacion y calibracion por proceso;
\item
  auditoria post-operacion y benchmarking de recursos;
\item
  mejora de ofertas y planificacion mediante distribuciones de duracion;
\item
  menor coste interno y mejor cumplimiento de SLA en la terminal de
  Vigo;
\item
  diferenciacion de servicios de ingenieria y operacion con evidencia
  digital.
\end{enumerate}

No se debe poner precio ni prometer ahorro antes de conocer volumen,
coste por hora, penalizaciones, demurrage, productividad y comprador
real.

\subsection{Preguntas de
descubrimiento}\label{preguntas-de-descubrimiento}

\begin{itemize}
\tightlist
\item
  ¿Cual de los tres productos tiene mas operaciones finalizadas y
  timestamps?
\item
  ¿Trace Port guarda cada evento o solo el informe final?
\item
  ¿Que constituye una desviacion economicamente material?
\item
  ¿Se almacena el plan original y cada replanificacion?
\item
  ¿Se pueden enlazar packing list, operacion, turno, recurso e
  incidencia?
\item
  ¿Que decisiones se pueden cambiar durante una operacion?
\item
  ¿Quien compraria el modulo: Kaleido interno, terminal, cargador o
  cliente final?
\item
  ¿Que coste tiene una hora de retraso, una espera y una falsa alarma?
\item
  ¿Se puede hacer una exportacion anonimizada de 5 operaciones esta
  semana?
\end{itemize}

\end{document}
